\documentclass[12pt,a4paper]{article}

% --- PAKIETY PODSTAWOWE I JĘZYK POLSKI ---
\usepackage[utf8]{inputenc}
\usepackage[T1]{fontenc}
\usepackage[polish]{babel}
\usepackage{lmodern}
\usepackage{geometry}
\geometry{
    top=2.5cm,
    bottom=2.5cm,
    left=2.5cm,
    right=2.5cm
}

% --- PAKIETY GRAFICZNE I TABELE ---
\usepackage{graphicx}
\usepackage{booktabs}
\usepackage{tabularx}
\usepackage{float}
\usepackage{xcolor}

% --- LISTINGI KODU ŹRÓDŁOWEGO ---
\usepackage{listings}
\definecolor{codegreen}{rgb}{0,0.6,0}
\definecolor{codegray}{rgb}{0.5,0.5,0.5}
\definecolor{codepurple}{rgb}{0.58,0,0.82}
\definecolor{backcolour}{rgb}{0.96,0.96,0.96}

\lstdefinestyle{mystyle}{
    backgroundcolor=\color{backcolour},   
    commentstyle=\color{codegreen},
    keywordstyle=\color{blue}\bfseries,
    numberstyle=\tiny\color{codegray},
    stringstyle=\color{codepurple},
    basicstyle=\ttfamily\footnotesize,
    breakatwhitespace=false,         
    breaklines=true,                 
    captionpos=b,                    
    keepspaces=true,                 
    numbers=left,                    
    numbersep=5pt,                  
    showspaces=false,                
    showstringspaces=false,
    showtabs=false,                  
    tabsize=2,
    frame=single,
    rulecolor=\color{lightgray}
}
\lstset{style=mystyle}

% --- ODNOŚNIKI HYPERREF ---
\usepackage[colorlinks=true, linkcolor=black, citecolor=blue, urlcolor=blue]{hyperref}

% ========================================================
% STRONA TYTUŁOWA
% ========================================================
\begin{document}

\begin{titlepage}
    \centering
    \vspace*{1cm}
    {\Large\textbf{POLITECHNIKA POZNAŃSKA}\par}
    {\large Wydział Informatyki i Telekomunikacji\par}
    {\large Instytut Informatyki – Zakład Systemów Informatycznych\par}
    
    \vspace{2.5cm}
    {\huge\bfseries DOKUMENTACJA PROJEKTOWA\par}
    \vspace{0.5cm}
    {\Large\textbf{Dedykowany System Zarządzania Treścią (CMS) \\ dla Restauracji Gastronomicznej}\par}
    \vspace{0.3cm}
    {\large Projekt realizowany w technologii Node.js, Express, EJS i SQLite\par}
    
    \vspace{3cm}
    \begin{minipage}{0.45\textwidth}
        \begin{flushleft} \large
            \textbf{Autorzy:}\\
            Imię i Nazwisko 1 (nr indeksu)\\
            Imię i Nazwisko 2 (nr indeksu)\\
            Kierunek: Informatyka\\
            Specjalność: ZTI (st. II stopnia)
        \end{flushleft}
    \end{minipage}
    \hfill
    \begin{minipage}{0.45\textwidth}
        \begin{flushright} \large
            \textbf{Prowadzący:}\\
            dr inż. Michał Apolinarski\\
            \vspace{0.5cm}
            \textbf{Przedmiot:}\\
            Systemy Zarządzania Treścią\\
            Laboratorium
        \end{flushright}
    \end{minipage}

    \vfill
    {\large Poznań, 2026\par}
\end{titlepage}

% ========================================================
% SPIS TREŚCI
% ========================================================
\tableofcontents
\newpage

% ========================================================
% 1. CHARAKTERYSTYKA OGÓLNA PROJEKTU
% ========================================================
\section{Charakterystyka ogólna projektu}
Celem projektu jest opracowanie i implementacja dedykowanego, wielofunkcyjnego systemu zarządzania treścią (Content Management System – CMS) dla restauracji gastronomicznej \textit{Sarab Restaurant}. 

System rozwiązuje problem dynamicznego zarządzania informacjami w lokalu gastronomicznym, takimi jak:
\begin{itemize}
    \item bieżąca karta menu, ceny oraz aktywne \textbf{Danie Dnia},
    \item obsługa rezerwacji stolików w czasie rzeczywistym,
    \item publikacja i moderacja opinii klientów,
    \item harmonogram godzin otwarcia lokalu,
    \item prezentacja zespołu kucharzy oraz kanałów kontaktu z personelem.
\end{itemize}

Aplikacja składa się z dwóch powiązanych części: portalu publicznego dla klientów oraz bezpiecznego panelu administracyjnego z autoryzacją opartą na rolach (Role-Based Access Control – RBAC).

% ========================================================
% 2. WYMAGANIA SYSTEMU
% ========================================================
\section{Wymagania systemu}

\subsection{Wymagania funkcjonalne z podziałem na aktorów}

W systemie zidentyfikowano trzech głównych aktorów: \textbf{Gościa restauracji (Klienta)}, \textbf{Pracownika} oraz \textbf{Właściciela}.

\begin{table}[H]
\centering
\small
\begin{tabularx}{\textwidth}{lX}
\toprule
\textbf{Aktor} & \textbf{Wymagania funkcjonalne} \\
\midrule
\textbf{Gość (Klient)} & 
\begin{itemize}
    \item Przeglądanie karty dań z podziałem na kategorie (Burgery, Pizza, Kurczak, Wrapy, Desery, Makarony).
    \item Podgląd szczegółów dania w oknie modalnym (kalorie, czas przygotowania, skład).
    \item Składanie rezerwacji stolika poprzez interaktywny formularz.
    \item Wystawianie ocen i opinii o lokalu (ocena w skali 1--5 gwiazdek).
    \item Wysyłanie zapytań przez formularz kontaktowy.
\end{itemize} \\
\midrule
\textbf{Pracownik} & 
\begin{itemize}
    \item Logowanie do panelu administracyjnego CMS z poziomu konta pracowniczego.
    \item Dodawanie, edycja i usuwanie potraw w karcie menu (pełny CRUD).
    \item Szybkie wyznaczanie i zmiana aktualnego \textbf{Dania Dnia} jednym kliknięciem.
    \item Przeglądanie nadesłanych rezerwacji stolików oraz zmiana ich statusu (\textit{Oczekująca}, \textit{Potwierdzona}, \textit{Odrzucona}, \textit{Zakończona}).
    \item Aktualizacja godzin otwarcia lokalu dla poszczególnych dni tygodnia.
    \item Przeglądanie i obsługa wiadomości z formularza kontaktowego.
\end{itemize} \\
\midrule
\textbf{Właściciel} & 
\begin{itemize}
    \item Posiadanie wszystkich uprawnień przypisanych do roli Pracownika.
    \item Moderacja i zatwierdzanie opinii klientów przed ich publikacją na stronie.
    \item Zarządzanie kontami użytkowników CMS (dodawanie nowych pracowników, usuwanie kont).
    \item Edycja profilu szefów kuchni prezentowanych w witrynie.
    \item Konfiguracja globalnych danych lokalu (telefon, e-mail, adres, hasło reklamowe).
\end{itemize} \\
\bottomrule
\end{tabularx}
\caption{Zestawienie wymagań funkcjonalnych z podziałem na aktorów}
\end{table}

\subsection{Wymagania pozafunkcjonalne}
\begin{enumerate}
    \item \textbf{Responsywność (RWD):} Zarówno portal klienta, jak i panel CMS muszą poprawnie skalować się na urządzeniach mobilnych, tabletach oraz desktopach (oparte na Bootstrap 5).
    \item \textbf{Bezpieczeństwo sesji:} Ochrona tras administracyjnych przed nieautoryzowanym dostępem za pomocą middleware weryfikującego stan sesji HTTP (\texttt{express-session}) oraz poziom uprawnień roli.
    \item \textbf{Wydajność bazy danych:} Zastosowanie synchronicznego i bezserwerowego sterownika \texttt{better-sqlite3}, eliminującego opóźnienia sieciowe związane z bazą.
    \item \textbf{Wielojęzyczność i skalowalność:} Przygotowanie struktury szablonów pod łatwą internacjonalizację (I18n).
\end{enumerate}

% ========================================================
% 3. ARCHITEKTURA SYSTEMU I TECHNOLOGIE
% ========================================================
\section{Architektura systemu i technologie}

\subsection{Środowisko i stos technologiczny}
Aplikacja została wykonana w architekturze wielowarstwowej zgodnie ze wzorcem \textbf{MVC (Model-View-Controller)}:
\begin{itemize}
    \item \textbf{Środowisko uruchomieniowe:} Node.js v22 (silnik JavaScript V8).
    \item \textbf{Framework aplikacji backendowej:} Express.js v5.
    \item \textbf{Silnik szablonów (Server-Side Rendering):} EJS (Embedded JavaScript).
    \item \textbf{System zarządzania bazą danych:} SQLite z modułem \texttt{better-sqlite3}.
    \item \textbf{Zarządzanie stanem i sesjami:} \texttt{express-session}.
    \item \textbf{Interfejs graficzny:} HTML5, CSS3, Bootstrap 5.3, Bootstrap Icons, FontAwesome 6, biblioteki animacji (AOS, Swiper).
\end{itemize}

\subsection{Struktura organizacji plików}
\begin{lstlisting}[language=bash, caption={Organizacja katalogów projektu}]
Systemy Zarzadzania Trescia/
|-- server.js               # Glowny kontroler aplikacji Express
|-- package.json            # Zaleznosci i skrypty npm
|-- database/               # MODEL: konfiguracja bazy, init.js i sqlite
|-- routes/                 # KONTROLER: adminRoutes.js (CMS, ochrona rbac)
|-- views/                  # WIDOKI (EJS):
|   |-- client/             # 7 podstron klienta (Home, Menu, About, itd.)
|   |   `-- partials/       # Wspolne elementy (header, footer, banner)
|   `-- admin/              # Widoki CMS (dashboard, menu, rezerwacje, users)
|-- client/                 # Statyczne assety klienta (CSS, JS, img)
|-- admin/                  # Statyczne assety panelu AdminLTE v4
`-- docs/                   # Dokumentacja projektu (sprawozdanie LaTeX)
\end{lstlisting}

% ========================================================
% 4. SCHEMAT BAZY DANYCH
% ========================================================
\section{Schemat bazy danych}

\subsection{Model relacyjny}
Baza danych została zrealizowana w oparciu o system relacyjny SQLite z włączoną integralnością kluczy obcych (\texttt{PRAGMA foreign\_keys = ON}).

Zdefiniowano następujące tabele:
\begin{enumerate}
    \item \textbf{users} (\underline{id}, username, password, full\_name, role, created\_at)
    \item \textbf{categories} (\underline{id}, name, slug, item\_count\_label, image)
    \item \textbf{menu\_items} (\underline{id}, category\_id $\rightarrow$ categories(id), name, description, price, old\_price, image, badge, calories, prep\_time, is\_dish\_of\_the\_day, is\_active)
    \item \textbf{reservations} (\underline{id}, full\_name, phone, email, guests, reservation\_date, reservation\_time, special\_requests, status, created\_at)
    \item \textbf{reviews} (\underline{id}, author\_name, role\_or\_city, rating, comment, is\_approved, created\_at)
    \item \textbf{chefs} (\underline{id}, name, role, experience, image, bio)
    \item \textbf{opening\_hours} (\underline{id}, day\_range, hours, is\_closed, sort\_order)
    \item \textbf{contact\_messages} (\underline{id}, name, email, phone, subject, message, is\_read, created\_at)
    \item \textbf{restaurant\_info} (\underline{key}, value)
\end{enumerate}

% ========================================================
% 5. DIAGRAMY UML
% ========================================================
\section{Diagramy UML}

\subsection{Diagram przypadków użycia (Use Case)}
\textit{[W tym miejscu wklej wygenerowany diagram przypadków użycia dla Klienta, Pracownika i Właściciela]}.

\subsection{Diagram klas / encji}
\textit{[W tym miejscu wklej diagram klas przedstawiający encje modeli bazodanowych oraz ich powiązania]}.

% ========================================================
% 6. NAJWAŻNIEJSZE METODY I FRAGMENTY KODU
% ========================================================
\section{Najważniejsze metody i fragmenty kodu aplikacji}

\subsection{Mechanizm autoryzacji i ochrony ról (RBAC)}
\begin{lstlisting}[language=JavaScript, caption={Middleware sprawdzający uprawnienia Właściciela w routes/adminRoutes.js}]
function requireOwner(req, res, next) {
    if (!req.session.user || req.session.user.role !== 'owner') {
        return res.status(403).send('Brak uprawnien. Sekcja dla Wlasciciela.');
    }
    next();
}
\end{lstlisting}

\subsection{Logika przełączania Dania Dnia}
\begin{lstlisting}[language=JavaScript, caption={Obsługa wyznaczania Dania Dnia z zerowaniem poprzedniego stanu}]
router.post('/menu/toggle-daily/:id', (req, res) => {
    const dish = db.prepare('SELECT is_dish_of_the_day FROM menu_items WHERE id = ?').get(req.params.id);
    if (dish) {
        if (dish.is_dish_of_the_day === 1) {
            db.prepare('UPDATE menu_items SET is_dish_of_the_day = 0 WHERE id = ?').run(req.params.id);
        } else {
            db.prepare('UPDATE menu_items SET is_dish_of_the_day = 0').run();
            db.prepare('UPDATE menu_items SET is_dish_of_the_day = 1 WHERE id = ?').run(req.params.id);
        }
    }
    res.redirect('/admin/menu');
});
\end{lstlisting}

% ========================================================
% 7. PODSUMOWANIE I PODZIAŁ PRAC
% ========================================================
\section{Podsumowanie}

\subsection{Podział prac w zespole}
\begin{itemize}
    \item \textbf{Autor 1:} Implementacja struktury aplikacji Express, pocięcie szablonów EJS, integracja bazy danych SQLite, wdrożenie panelu CMS AdminLTE.
    \item \textbf{Autor 2:} Opracowanie dokumentacji projektowej LaTeX, przygotowanie diagramów UML i ERD, testy modułowe i formularzy.
\end{itemize}

\subsection{Perspektywa rozwoju}
\begin{itemize}
    \item Wdrożenie fizycznego uploadu plików ze skalowaniem zdjęć przy użyciu biblioteki \texttt{multer} i \texttt{sharp}.
    \item Rozszerzenie modułu zamówień online o obsługę koszyka zakupowego i płatności internetowych.
\end{itemize}

\end{document}
